\documentclass{article}
\usepackage{graphicx} % Required for inserting images
\usepackage{amsmath}
\usepackage{amssymb}
\usepackage[margin=1in]{geometry}
\usepackage[most]{tcolorbox}
\usepackage{bm}
\usepackage{lipsum} % for dummy text

\newcommand{\mydivider}{\vspace{1em}\hrule\vspace{1em}}

\tcbset{
    enhanced,
    breakable,
    skin first=enhanced,
    skin middle=enhanced,
    skin last=enhanced,
}

%%% BEGIN MACROS
% type your macros here
%%% END MACROS

\title{CS 577 --- Deep Learning --- Homework 1}
\author{}
\date{}

\begin{document}

\maketitle
\textbf{Read these instructions carefully}:
\begin{itemize}
  \item
In the \LaTeX~source code, type your answer in between
``\verb|%%% BEGIN ANSWER|'' and
``\verb|%%% END ANSWER|''.
        For advanced \LaTeX users,
        you can use your custom macros if you wish
        by placing them
        between
``\verb|%%% BEGIN MACROS|'' and
``\verb|%%% END MACROS|'' in the header.
        Do not modify anything else.


  \item
There is a PDF file \verb|latex_symbols.pdf| which lists \LaTeX~symbols that you will need later.


  \item
The first section ``Introduction to Latex (not graded)'' is a tutorial. \textbf{It is not graded}. Feel free to skip it.

  \item The second section titled ``\LaTeX~Problems'' \textbf{will be graded}. You will recreate content from the slides in Lecture 1. You will receive full credit as long as your answer is sufficiently similar visually.

  \item The third section ``Perceptron Problem'' \textbf{will be graded}.

  \item When typing mathematics, you can use either the inline mode or display mode (see below). It is up to you. Please make the choice that results in good readability.

  \item Turn in both your \verb|.tex| file and the generated \verb|.pdf| file.
\end{itemize}



\section{Introduction to Latex (not graded)}


\textbf{[0] points --- inline mode:}
Mathematics goes inside a pair of single dollar sign like this: ``\verb|$...$|''.
For example,
``\verb|$\pi \approx 3.14$|'' renders as
``$\pi \approx 3.14$''

 Write the formula for the circumference $C$ of the circle with radius $r$

\begin{tcolorbox}{}
\textbf{Answer: }

%%% BEGIN ANSWER
% type your answer here
%%% END ANSWER


\end{tcolorbox}


\mydivider
\noindent\textbf{[0] points --- display mode:} Mathematics goes inside a pair of double dollar signs for display mode like this: ``\verb|$$...$$|''. For example,
``\verb|$$\sin(\theta)^2 + \cos(\theta)^2 = 1$$|'' renders as

$$\sin(\theta)^2 + \cos(\theta)^2 = 1$$

 Type in display mode a function $f$ of $x$ and $y$ defined (you should use colon ``:'' followed by an equal sign, i.e., $:=$, for function definition) by $\sin(x)^2$ plus $\cos(y)^2$.


\begin{tcolorbox}{}
\textbf{Answer: }

%%% BEGIN ANSWER
% type your answer here
%%% END ANSWER

\end{tcolorbox}




\mydivider
\noindent\textbf{[0] points --- superscript:} You saw that the caret symbol \verb|^| is for superscript. Now, try taking $x$ to the $1/3$ power by directly typing \verb|x^| followed by \verb|1/3|...

\begin{tcolorbox}{}
\textbf{Answer: }

%%% BEGIN ANSWER
% type your answer here
%%% END ANSWER

\end{tcolorbox}

If you did this exactly, you probably noticed that it's not quite right. Let's fix that by surrounding the \verb|1/3| with \verb|{...}|.

\begin{tcolorbox}{}
\textbf{Answer: }

%%% BEGIN ANSWER
% type your answer here
%%% END ANSWER

\end{tcolorbox}



\mydivider

\noindent\textbf{[0] points --- subscript:} Subscripts in LaTeX are created using the underscore symbol \verb|_|. For example, typing \verb|x_i| will produce $x_i$. Now, try typing $x$ with the subscript ``$100$''.

\begin{tcolorbox}{}
\textbf{Answer: }

%%% BEGIN ANSWER
% type your answer here
%%% END ANSWER

\end{tcolorbox}


\mydivider

\noindent\textbf{[0] points --- summation notation:} The command \verb|\sum| is used to produce the summation symbol in LaTeX. For example, \verb|\sum_{...}^{...} | produces $\sum_{...}^{...}$. Now, write the sum of $1$ squared, $2$ squared, $\ldots$, $n$ squared using summation notation.

\begin{tcolorbox}{}
\textbf{Answer: }

%%% BEGIN ANSWER
% type your answer here
%%% END ANSWER

\end{tcolorbox}

\mydivider


\noindent\textbf{[0] points --- blackboard font:} The set of real numbers is typically written using the blackboard font in LaTeX. To use this font, surround the symbol with \verb|\mathbb{...}|. Now, write the set of real numbers (capital $R$ in blackboard font).

\begin{tcolorbox}{}
\textbf{Answer: }

%%% BEGIN ANSWER
% type your answer here
%%% END ANSWER

\end{tcolorbox}

\mydivider

\noindent\textbf{[0] points --- minimization notation:} Minimization in LaTeX is similar to summation, except there is no superscript. Now, write the notation for ``minimization of $x$ squared where $x$ is in the set of the reals''. The symbol for ``belongs'' is \verb|\in|.

\begin{tcolorbox}{}
\textbf{Answer: }

%%% BEGIN ANSWER
% type your answer here
%%% END ANSWER

\end{tcolorbox}

\mydivider

\noindent\textbf{[0] points --- bold symbols:} To type a bold symbol in LaTeX, use the command \verb|\mathbf{...}|, where the \verb|...| is replaced with your symbol.  Now, write $X$  (uppercase) in bold.

\begin{tcolorbox}{}
\textbf{Answer: }

%%% BEGIN ANSWER
% type your answer here
%%% END ANSWER

\end{tcolorbox}

For bolding Greek symbols, you need to use \verb|\bm{...}| instead of \verb|\mathbf{...}|.  Write ``theta'' in bold.

\begin{tcolorbox}{}
\textbf{Answer: }

%%% BEGIN ANSWER
% type your answer here
%%% END ANSWER

\end{tcolorbox}

\mydivider

\noindent\textbf{[0] points --- uppercase Greek letters:} To type an uppercase Greek letter in LaTeX, simply write the name of the letter with an initial capital letter, like \verb|\Theta| for $\Theta$. Note that not all Greek letters have an uppercase version in LaTeX (e.g., there is no capital "alpha"). Now, type the uppercase version of delta.

\begin{tcolorbox}{}
\textbf{Answer: }

%%% BEGIN ANSWER
% type your answer here
%%% END ANSWER

\end{tcolorbox}


\mydivider

\noindent\textbf{[0] points --- gradient symbol:} Other commands or symbols in LaTeX can also be subscripted or superscripted. The symbol for the gradient is \verb|\nabla|. Now, type the symbol for the partial derivative with respect to bold $\mathbf{x}$ (just the "nabla" and the "x" part).

\begin{tcolorbox}{}
\textbf{Answer: }

%%% BEGIN ANSWER
% type your answer here
%%% END ANSWER

\end{tcolorbox}


\mydivider

\noindent\textbf{[0] points --- fractions:} Fractions in LaTeX are written using the command \verb|\frac{numerator}{denominator}|. Now, write Newton's law of universal gravitation:
$F$ is the force of gravity, $G$ is the gravitational constant, $m_1$ and $m_2$ are the masses of the objects, and $r$ is the distance between the objects.

\begin{tcolorbox}{}
\textbf{Answer: }

%%% BEGIN ANSWER
% type your answer here
%%% END ANSWER

\end{tcolorbox}


\subsection{Linear algebra}

\noindent\textbf{[0] points --- matrices:} You can write matrices in LaTeX using the \verb|bmatrix| environment. For example, the matrix
\[
\begin{bmatrix}
a & b \\
c & d
\end{bmatrix}
\]
is written as:

\begin{verbatim}
\begin{bmatrix}
a & b \\
c & d
\end{bmatrix}
\end{verbatim}

Now, write down the formula for the inverse of this matrix.

\begin{tcolorbox}{}
\textbf{Answer: }

%%% BEGIN ANSWER
% type your answer here
%%% END ANSWER

\end{tcolorbox}



\mydivider

\noindent\textbf{[0] points --- big matrices:} You can write a big matrix in LaTeX using dots for representing a sequence of elements. Here's an example of a large matrix:

\[
\begin{bmatrix}
x_{11} & \cdots & x_{1N} \\
\vdots & \ddots & \vdots \\
x_{N1} & \cdots & x_{NN}
\end{bmatrix}
\]

In this example, \verb|\cdots| represents horizontal (centered) dots, \verb|\vdots| represents vertical dots, and \verb|\ddots| represents diagonal dots.

Now, create a column vector with $x_1, \ldots, x_N$ and a row vector with $x_1, \ldots, x_N$. Hint: You only need \verb|&| but not \verb|\\| for the row vector.

\begin{tcolorbox}{}
\textbf{Answer: }

%%% BEGIN ANSWER
% type your answer here
%%% END ANSWER

\end{tcolorbox}


\subsection{Lists}


\noindent\textbf{[0] points --- nested lists:} In LaTeX, you can create a nested bulleted list using the \verb|itemize| environment like this:

\begin{verbatim}
\begin{itemize}
    \item A
    \item B
    \begin{itemize}
        \item X
        \item Y
        \item Z
    \end{itemize}
    \item C
\end{itemize}
\end{verbatim}

This will produce the following nested list:

\begin{itemize}
    \item A
    \item B
    \begin{itemize}
        \item X
        \item Y
        \item Z
    \end{itemize}
    \item C
\end{itemize}

To create a numbered list, use the \verb|enumerate| environment instead of \verb|itemize|.
Nested lists can be a hacky way to write pseudo algorithms (with code indentation) in LaTeX. Now, write the pseudocode to compute the $n$-th Fibonacci number using nested numbered list.

\begin{tcolorbox}{}
\textbf{Answer: }

%%% BEGIN ANSWER
% type your answer here
%%% END ANSWER

\end{tcolorbox}




\section{\LaTeX~Problems}

Note: The following problems refer to slides from Lecture 01.

\vspace{1em}

\noindent\textbf{[2] points --- Empirical risk minimization:} Write the empirical risk minimization (ERM)
on slide number ``27/73''  (ignore my handwritten notes) of . Consult the \verb|latex_symbols.pdf| page 4 for the matrix transpose symbol.

\begin{tcolorbox}{}
\textbf{Answer: }

%%% BEGIN ANSWER
% type your answer here
%%% END ANSWER

\end{tcolorbox}

\mydivider

\noindent\textbf{[2] points --- Gradient:} Write the gradient
on slide number ``29/73''  (ignore my handwritten notes).
Consult the
\verb|latex_symbols.pdf| page 4 for the partial derivative symbol.
\begin{tcolorbox}{}
\textbf{Answer: }

%%% BEGIN ANSWER
% type your answer here
%%% END ANSWER

\end{tcolorbox}

\mydivider


\noindent\textbf{[3] points --- Stochastic gradient descent algorithm:} Write the
content on slide number ``42/73''  (ignore my handwritten notes) starting from ``Let...''. (Ignore the slide numbers in your answer.)
The curly brackets are \verb|\{...\}|
and the left arrow is \verb|\gets|.
\begin{tcolorbox}{}
\textbf{Answer: }

%%% BEGIN ANSWER
% type your answer here
%%% END ANSWER

\end{tcolorbox}


\mydivider


\noindent\textbf{[3] points --- Perceptron algorithm:} Write the
content inside the ``Perceptron update'' blue box on slide number ``53/73''  (ignore my handwritten notes).
(Just the pseudocode. Ignore the title text of the box.)

\begin{tcolorbox}{}
\textbf{Answer: }

%%% BEGIN ANSWER
% type your answer here
%%% END ANSWER

\end{tcolorbox}




\section{Perceptron Problem}
\noindent\textbf{[10] points --- Perceptron calculation:} In this problem, you will run the Perceptron update for the following dataset for \(T = 8\) iterations. It is possible for the Perceptron to converge before all \(T\) iterations. In that case, you may declare that ``the perceptron converges after this iteration'' and stop.
\begin{itemize}
  \item $\mathbf{x}^{(1)} = \begin{bmatrix} 1 \\ 0 \end{bmatrix}$ \quad ``east'',
        \quad $y^{(1)} = +1$
    \item $\mathbf{x}^{(2)} = \begin{bmatrix} 1 \\ 1 \end{bmatrix}$ \quad ``northeast'',
        \quad $y^{(2)} = +1$
    \item $\mathbf{x}^{(3)} = \begin{bmatrix} 0 \\ 2 \end{bmatrix}$ \quad ``north'',
        \quad $y^{(3)} = -1$
    \item $\mathbf{x}^{(4)} = \begin{bmatrix} -1 \\ 1 \end{bmatrix}$ \quad ``northwest'',
        \quad $y^{(4)} = -1$
    \item $\mathbf{x}^{(5)} = \begin{bmatrix} -1 \\ 0 \end{bmatrix}$ \quad ``west'',
        \quad $y^{(5)} = -1$
    \item $\mathbf{x}^{(6)} = \begin{bmatrix} -1 \\ -1 \end{bmatrix}$ \quad ``southwest'',
        \quad $y^{(6)} = -1$
    \item $\mathbf{x}^{(7)} = \begin{bmatrix} 0 \\ -1 \end{bmatrix}$ \quad ``south'',
        \quad $y^{(7)} = +1$
    \item $\mathbf{x}^{(8)} = \begin{bmatrix} 1 \\ -1 \end{bmatrix}$ \quad ``southeast'',
        \quad $y^{(8)} = +1$
\end{itemize}
Write your answer in the box below. Show all your calculations.
Go through each iteration thoroughly.
Make sure that you address
\begin{itemize}
  \item
What is $\mathbf{w}$ at the beginning/end of an iteration?
  \item What is the calculation that goes into checking the ``If/Else'' statement?

\end{itemize}

\begin{tcolorbox}{}
\textbf{Answer: }

%%% BEGIN ANSWER
% type your answer here
%%% END ANSWER

\end{tcolorbox}
\end{document}
